\section{Further questions and verification}
The theorem resolves the nondivisible cardinalities in an explicit
sufficient range. Determining the least threshold for the balanced
four-block construction remains open. In higher codimension, a related
question is whether the extremizers at nondivisible cardinalities arise
by deforming replicated maximal equiangular tight frames, when such
frames exist. An extension would require both a localization theorem
for the corresponding projective configuration and a strict comparison
for its possible integer multiplicities.

Two features of the present argument may be useful for these questions.
The common deficit budget couples radial and angular deviations during
localization. With the candidate stress fixed, the kernel functional
$\Gamma_K$ is exactly quadratic under positive radial rescalings of
the columns. The radial--tangential decomposition in the local
comparison uses this property to control the remainder. These features
reduce the loss in passing from a global bound to a finite neighborhood
of the candidate.

The proof uses computer assistance to verify the finite inequalities in
Sections~\ref{app:candidate}--\ref{app:certificates}. The supplementary
package provides rational certificates, an arithmetic verifier that checks
them without repeating the certificate search, and instructions for
reproducing the verification.
Section~\ref{app:certificates} specifies the verifier's scope and the
analytic arguments that extend the finite checks to the stated range.
